%%%%%%%%%%  scan idx 383  |  folio 370  %%%%%%%%%%
\SGAPrintPage{370}%
and a sheaf \(\underline{F}_{\lambda'}\) of finite presentation on \(X_{\lambda'}\), such that \(\underline{F}=u'^*_{\lambda'}(\underline{F}_{\lambda'})\) (\(u'_{\lambda'}\) being the morphism \(X\to X_{\lambda'}\)). Since the \(X_\mu\) and the \(\underline{F}_\mu\) are of finite presentation, we can, by EGA IV 11.2.6, find an index \(\lambda''\) such that \(\underline{F}_{\lambda''}\) is flat over \(S_{\lambda''}\). All the hypotheses of 1.3 are therefore satisfied, which completes the proof.

\sgaStableTarget{sga6.e06.cor.1-4}\textbf{Corollary 1.4.} \uline{Under the hypotheses of} 1.3, \uline{there exists a graded} \(\Sym(\underline{E})\)\uline{-Module} \(\underline{M}\), \uline{of finite presentation}, \uline{and flat over} \(\underline{O}_S\), \uline{such that} \(\underline{F}=\Projz(\underline{M})\).

Let \(n_0\) be an integer such that the properties of 1.3 hold. If \(\underline{M}=\bigoplus_{n\geq n_0} f_*(\underline{F}(n))\), we have seen that \(\underline{M}\) is of finite presentation, and it is flat over \(S\), since the \(f_*(\underline{F}(n))\), for \(n\geq n_0\), are locally free on \(S\).

Moreover, \(\underline{F}=\Projz(\underline{M})\) by EGA II 3.4.

Let us turn to part b) of the proof.

\sgaStableTarget{sga6.e06.lem.1-5-1}\textbf{Lemma 1.5.1.} \uline{Let} \(A\) \uline{be a ring}, \(B\) \uline{a positively graded} \(A\)-\uline{algebra} \uline{such that} \(B_0=A\), \(M\) \uline{a graded} \(B\)-\uline{module}. \uline{We suppose that} \(M\) \uline{is projective of finite type}, \uline{as an ungraded} \(B\)-\uline{module}. \uline{Let} \(\widetilde B\) \uline{be the sheaf of algebras on} \(\Spec(A)\) \uline{defined by} \(B\), \(\widetilde M\) \uline{the} \(\widetilde B\)-\uline{Module associated with} \(M\); \uline{then} \(\widetilde M\) \uline{is, locally on} \(\Spec(A)\), \uline{free of finite type over} \(\widetilde B\), \uline{as a graded} \(\widetilde B\)-\uline{Module} (1.2).

Since \(M\) is of finite type, its degrees are bounded below. Moreover, \(M\otimes_B A\) is projective of finite type over \(A\), and consequently has only finitely many nonzero homogeneous components, themselves projective of finite type over \(A\); by localizing on \(\Spec(A)\), we can therefore suppose that \(M\otimes_B A\) is free of finite type over \(A\). Finally, since \(M\) is projective, \(\Tor_1^B(M,A)=0\). These properties imply that \(M\) is free of finite type over \(B\) as a graded module (see Bourbaki, Alg. ch. II §11 prop. 7), so \(\widetilde M\) is free of finite type over \(\widetilde B\), q.e.d.

%%%%%%%%%%  scan idx 384  |  folio 371  %%%%%%%%%%
\SGAPrintPage{371}%
\sgaStableTarget{sga6.e06.lem.1-5-2}\textbf{Lemma 1.5.2.} \uline{Let} \(A\) \uline{be a ring}, \(B\) \uline{an} \(A\)-\uline{algebra of finite presentation}, \(M\) \uline{a} \(B\)-\uline{module of finite presentation}. \uline{We suppose} \(B\) \uline{and} \(M\) \uline{flat over} \(A\), \uline{and} \uline{that for every residue field} \(k\) \uline{of} \(A\), \(M\otimes_A k\) \uline{has projective dimension} \(\leq n\) \uline{over} \(B\otimes_A k\) (\uline{which is the case if} \(B=A[t_0,\ldots,t_r]\)). \uline{Under these conditions}, \(M\) \uline{admits a resolution of length} \(n\) \uline{by projective} \(B\)-\uline{modules of finite type}:
\[
  0\longrightarrow L_n\longrightarrow\ldots\longrightarrow L_0\longrightarrow M\longrightarrow 0.
\]
This lemma follows immediately from EGA IV 12.3.2.

\sgaStableTarget{sga6.e06.prop.1-6}\textbf{Proposition 1.6.} \uline{Let} \(S\) \uline{be a quasi-compact scheme}, \(\underline{E}\) \uline{a} \(\underline{O}_S\)\uline{-Module locally free of rank} \(r+1\), \(\underline{B}\) \uline{its symmetric algebra}, \uline{with its usual grading}, \(\underline{M}\) \uline{a} \(\underline{B}\)\uline{-Module that is graded and of finite presentation}, \uline{flat over} \(S\). \uline{Then there exists a graded resolution}:
\[
  0\longrightarrow \underline{L}_{r+1}\longrightarrow\ldots\longrightarrow \underline{L}_0\longrightarrow \underline{M}\longrightarrow 0,
\]
\uline{where the} \(\underline{L}_i\) \uline{are} \(\underline{B}\)\uline{-Modules that are graded and locally free of finite type} (\uline{as graded} \(\underline{B}\)\uline{-Modules}).

Let us first show that for every graded \(\underline{B}\)-Module \(\underline{M}\) satisfying the conditions of the proposition, there exists a surjective graded homomorphism: \(\underline{L}\to \underline{M}\), where \(\underline{L}\) is a \(\underline{B}\)-Module locally free of finite type.

Since \(\underline{M}\) is of finite type, there exists an integer \(N\) such that \(\underline{M}\) is generated over \(\underline{B}\) by its components of degree \(\leq N\). We then put:
\[
  \underline{L}=\bigoplus_{i\leq N} \underline{M}_i\otimes_{\underline{O}_S}\underline{B}(-i);
\]
the homomorphism from \(\underline{L}\) to \(\underline{M}\) is defined using (1.2.1), and it is clear that it is surjective.

By 1.5.2, \(\underline{M}\) admits an (ungraded) resolution by \(\underline{B}\)-Modules locally free of finite type on every affine open set. Consequently, if
%%%%%%%%%%  scan idx 385  |  folio 372  %%%%%%%%%%
\SGAPrintPage{372}%
\(\underline{M}\) has a global partial resolution:
\[
  \underline{L}_i\longrightarrow \underline{L}_{i-1}\longrightarrow\ldots\longrightarrow \underline{L}_0\longrightarrow \underline{M}\longrightarrow 0,
\]
by \(\underline{B}\)-Modules locally free of finite type, the kernel \(\underline{Z}_i\) of \(\underline{L}_i\to \underline{L}_{i-1}\) is again a \(\underline{B}\)-Module of finite type (I. 2.11). It is even of finite presentation (by making the same argument with \(i+1\) instead of \(i\)). Finally, it is flat over \(S\). It therefore satisfies the same hypotheses as \(\underline{M}\), which allows us to construct step by step a global graded resolution of \(\underline{M}\) by graded \(\underline{B}\)-Modules locally free of finite type, by the initial remark. For \(i\geq r+1\), \(\underline{Z}_i\) is a projective \(\underline{B}\)-Module on every affine open set (1.5.2), and consequently is locally free as a graded \(\underline{B}\)-Module, by 1.5.1. We can then, in the graded resolution constructed, replace \(\underline{L}_{r+1}\) by \(\underline{Z}_r\), and the following terms by \(0\), whence the asserted resolution.

Let us finally prove point c) of the proof.

\sgaStableTarget{sga6.e06.prop.1-7}\textbf{Proposition 1.7.} \uline{Let} \(S\) \uline{be a quasi-compact scheme}, \(\underline{E}\) \uline{a} \(\underline{O}_S\)\uline{-Module locally free of finite type}, \(\underline{B}=\Sym(\underline{E})\), \uline{and} \(\underline{L}\) \uline{a graded} \(\underline{B}\)\uline{-Module}, \uline{locally free of finite type} (1.2). \uline{Then there exists a decreasing filtration of} \(\underline{L}\) \uline{by graded sub-}\(\underline{B}\)\uline{-Modules} \(\underline{L}^{(m)}\), \uline{such that}:
\[
\begin{array}{ll}
\text{i)}& \underline{L}^{(m)}=\underline{L}\quad\text{\uline{for} }m\leq m_0;\\[1mm]
\text{ii)}& \underline{L}^{(m)}=0\quad\text{\uline{for} }m\geq m_1;\\[1mm]
\text{iii)}& \text{\uline{for every} }m,\quad \operatorname{gr}^m(\underline{L})=\underline{L}^{(m)}/\underline{L}^{(m+1)}\\
&\quad \text{\uline{is isomorphic as a graded} }\underline{B}\text{\uline{-Module to a graded}}\\
&\quad \underline{B}\text{\uline{-Module of the form} }\underline{N}_m\otimes_{\underline{O}_S}\underline{B}(m),\\
&\quad \text{\uline{where} }\underline{N}_m\text{ \uline{is an} }\underline{O}_S\text{\uline{-Module locally free of finite type}.}
\end{array}
\]

Since \(S\) is quasi-compact, we can find integers \(n_0\) and \(n_1\) such that in a neighborhood of every point of \(S\), the elements of a homogeneous basis
%%%%%%%%%%  scan idx 386  |  folio 373  %%%%%%%%%%
\SGAPrintPage{373}%
of \(\underline{L}\) have degrees between \(n_0\) and \(n_1\). We then proceed by induction on \(N=n_1-n_0\).

When \(N=0\), i.e. \(n_1=n_0\), let \(n\) be the common value. The canonical homomorphism (1.2.1):
\[
  \underline{M}_n\otimes_{\underline{O}_S}\underline{B}(-n)\longrightarrow \underline{M}
\]
is then an isomorphism. Indeed, the verification is local, and by (1.2.3), we can then suppose that \(\underline{M}\simeq(\underline{B}(-n))^k\), and the result is trivial.

In the case \(N>0\), we again consider the homomorphism (1.2.1):
\[
  \underline{M}_{n_0}\otimes_{\underline{O}_S}\underline{B}(-n_0)\longrightarrow \underline{M}.
\]
We verify locally that this homomorphism is injective, and that its cokernel \(\underline{M}'\) is locally free of finite type, with local generators of degree between \(n_0+1\) and \(n_1\): indeed, in a neighborhood of every point of \(S\), \(\underline{M}_{n_0}\) is either zero or the nonzero homogeneous component of \(\underline{M}\) of lowest degree; using a decomposition of \(\underline{M}\) into direct summands of the form (1.2.2), we see in the latter case, using (1.2.3), that \(\underline{M}_{n_0}\otimes_{\underline{O}_S}\underline{B}(-n_0)\) is isomorphic to the direct summand \((\underline{B}(-n_0))^{k_{n_0}}\) of \(\underline{M}\), whence the asserted properties. We then apply the induction hypothesis to \(\underline{M}'\), whence the desired conclusion follows at once.

\sgaStableTarget{sga6.e06.sec.1.8}1.8. Returning to the proof of Theorem 1.1, let us show that \(K^{\bullet}(X)\) is generated as a \(K^{\bullet}(S)\)-module by the powers of \(\xi\).

If \(\underline{F}\) is a \(\underline{O}_X\)-Module locally free of finite type, there exists, by 1.4, a graded \(\underline{B}\)-Module of finite presentation \(\underline{M}\), flat over \(S\), such that \(\underline{F}=\Projz(\underline{M})\). By 1.6, \(\underline{M}\) admits a graded resolution by graded \(\underline{B}\)-Modules locally free of finite type. We deduce an exact sequence:
%%%%%%%%%%  scan idx 387  |  folio 374  %%%%%%%%%%
\SGAPrintPage{374}%
\[
  0\longrightarrow \Projz(\underline{L}_{r+1})\longrightarrow\ldots\longrightarrow
  \Projz(\underline{L}_0)\longrightarrow\Projz(\underline{M})\longrightarrow 0,
\]
and the following equality in \(K^{\bullet}(X)\):
\[
  \cl(\underline{F})=\sum_i(-1)^i\cl\bigl(\Projz(\underline{L}_i)\bigr).
\]
Using the filtration on the \(\underline{L}_i\) defined in 1.7, we find:
\[
  \cl(\Projz(\underline{L}_i))=\sum_j \cl\bigl(\Projz(\underline{N}_{ij}\otimes_{\underline{O}_S}\underline{B}(j))\bigr)
\]
Now \(\Projz(\underline{N}_{ij}\otimes_{\underline{O}_S}\underline{B}(j))=f^*(\underline{N}_{ij})(j)\); we therefore find:
\[
\begin{aligned}
  \cl(\Projz(\underline{L}_i))
  &=\sum_j \cl\bigl(f^*(\underline{N}_{ij})\bigr).\xi^j \\
  &=\sum_j f^*\bigl(\cl(\underline{N}_{ij})\bigr).\xi^j.
\end{aligned}
\]
We have thus finally expressed \(\cl(\underline{F})\) as a linear combination, with coefficients in \(K^{\bullet}(S)\), of the powers of \(\xi\).

\sgaStableTarget{sga6.e06.sec.1.9}1.9. To see that \(1,\xi,\ldots,\xi^{r+1}\) are linearly dependent over \(K^{\bullet}(S)\) we consider the canonical surjective homomorphism: \(f^*(\underline{E})\to\underline{O}_X(1)\), whose kernel \(\underline{F}\) is locally free of rank \(r\). We can therefore write:
\[
  \lambda^{r+1}\bigl(f^*(E)-\xi\bigr)=0,
  \qquad \text{with } E=\cl\underline{E},
\]
the operations \(\lambda^i\) on \(K^{\bullet}(X)\) being defined in V 2.2. We therefore obtain:
\[
  \sum_{i=0}^{r+1}\lambda^i(f^*(E))\lambda^{r+1-i}(-\xi)=0.
\]
Since \(\underline{O}_X(1)\) is an invertible sheaf, \(\lambda_t(\xi)=1+\xi t\), and consequently, \(\lambda_t(-\xi)=1/(1+\xi t)\), whence finally the relation:
\sgaStableTarget{sga6.e06.eq.1-9-1}\begin{equation}\tag{1.9.1}
  \sum_{i=0}^{r+1}(-1)^{r+1-i}\xi^{r+1-i}\lambda^i(E)=0.
\end{equation}

\sgaStableTarget{sga6.e06.sec.1.10}Let us finally show that \(1,\xi,\ldots,\xi^r\) are independent over \(K^{\bullet}(S)\)\SGAAuthorFootnote{*}{This proof is due to J.P. Jouanolou.}.

